% ============================================================
% paper/sections/04_axes_dof.tex
% Axes and Degrees of Freedom (ETS Projection) (finalized)
% ============================================================

\section{Axes and Degrees of Freedom}

\subsection{Role of Axes in the Enterprise Continuum}

Axes in $A_{\mathrm{EA}}$ define the admissible degrees of freedom of the
enterprise continuum.
They do not represent objectives, performance targets, maturity levels,
or normative dimensions.
Their sole function is to parameterize variation within the admissible
state space $\Omega_{\mathrm{EA}}$.

All axes described in this section are projected directly from
\texttt{ETS.K\_EA.v1.1}.
No additional axes are introduced, and none are reinterpreted beyond
their structural role as coordinates of admissible variation.

\subsection{Structural Axes}

Structural axes characterize the internal arrangement and coupling of
enterprise components.

\begin{itemize}
  \item \textbf{Coupling}: degree of dependency among components and
        subsystems.
  \item \textbf{Modularity}: extent to which components admit independent
        variation.
  \item \textbf{Centralization}: concentration of decision or control
        authority.
  \item \textbf{Redundancy}: presence of parallel components or pathways
        supporting persistence.
\end{itemize}

These axes determine the geometry and connectivity of
$\Omega_{\mathrm{EA}}$.
They do not encode preferred configurations or architectural ideals.

\subsection{Governance Axes}

Governance axes describe structural properties of rule formation,
decision propagation, and enforcement.

\begin{itemize}
  \item \textbf{Decision latency}: delay between signal detection and
        binding decision.
  \item \textbf{Rule enforceability}: degree to which declared rules are
        effectively applied.
  \item \textbf{Escalation integrity}: coherence of escalation paths
        under stress conditions.
\end{itemize}

These axes constrain admissible dynamics by shaping flows $J_{\mathrm{EA}}$
and the closure properties of cycles $C_{\mathrm{EA}}$.

\subsection{Operational Axes}

Operational axes parameterize execution-level properties of the
enterprise.

\begin{itemize}
  \item \textbf{Delivery coherence}: consistency between declared intent
        and realized outcomes.
  \item \textbf{Incident recovery closure}: ability to close recovery
        cycles after disruption.
  \item \textbf{Change lead time}: delay between decision and effective
        operational change.
\end{itemize}

Operational axes primarily affect persistence and stability rather than
existence.
Their violation may degrade continuity without immediately terminating
admissibility.

\subsection{Resource Axes}

Resource axes describe constraints associated with material, financial,
and human resources.

\begin{itemize}
  \item \textbf{Budget flexibility}: capacity to reallocate financial
        resources across commitments.
  \item \textbf{Skills coverage}: availability of competencies required
        to sustain cycles.
  \item \textbf{Supplier dependency}: degree of reliance on external
        providers.
\end{itemize}

These axes constrain potentials $P_{\mathrm{EA}}$ and contribute to
structural tension relative to thresholds $\Theta_{\mathrm{EA}}$.

\subsection{Information Axes}

Information axes characterize observability, trace construction, and
semantic consistency.

\begin{itemize}
  \item \textbf{Observability}: availability of reliable signals about
        internal state.
  \item \textbf{Traceability}: ability to reconstruct causal paths
        across decisions and actions.
  \item \textbf{Semantic alignment}: consistency of meaning across
        organizational units.
\end{itemize}

Loss of admissible observability directly constrains diagnosis and may
trigger termination conditions independent of other axes.

\subsection{Risk Axes}

Risk axes parameterize the accumulation and visibility of unmanaged
structural risk.

\begin{itemize}
  \item \textbf{Risk visibility}: ability to detect emerging structural
        risks.
  \item \textbf{Control gap rate}: rate at which unmitigated control
        deficits accumulate.
\end{itemize}

These axes interact with threshold components associated with inertia
and collapse.
They do not represent risk appetite or tolerance.

\subsection{Axes as Constraints, Not Goals}

Across all groups, axes function exclusively as constraints on
admissible variation.
No axis defines a direction of improvement, a target state, or a notion
of optimality.

An axis becomes diagnostically relevant only through its interaction
with thresholds $\Theta_{\mathrm{EA}}$, cycles $C_{\mathrm{EA}}$, and the
continuity measure $k_{\mathrm{EA}}$.
